% §4.6 What Comes Next
%
% Monograph register: Part I summary and Part II preview, compact.
% Cross-refs sec:axiom-1, sec:axiom-2, thm:kronecker, def:bernoulli-bool,
% ch:approximate-sets. DEFINED: sec:what-comes-next.

\section{What Comes Next}
\label{sec:what-comes-next}

Part~I is complete. The two axioms are stated (\Cref{sec:axiom-1},
\Cref{sec:axiom-2}), the Kronecker theorem shows what they license
(\Cref{thm:kronecker}), and \Cref{def:bernoulli-bool} gives
$\Bernoulli{\BoolType}$ a rigorous identity, of which every example in chapters
1 through 3 is an instance. Three of the framework's organizing principles are
now in view, each with its scope: the two axioms; closure of binary channels
under serial composition and Kronecker factorization of parallel queries, but
not closure of the two-rate type under Boolean formulas, where the exact
description is per latent assignment (\Cref{sec:boolean-formulas}); and the
Boolean thread that ran from the noisy bit to the formal definition without
leaving $\{\top, \bot\}$. Two more arrive in Part~II:
parametric parsimony\index{parametric parsimony}, already counted in
\Cref{cor:parameter-reduction} and priced in bits by the entropy bounds of
chapter~8, and space-accuracy duality\index{space--accuracy tradeoff}, the trade-off between a structure's size and its
error rates, which is Part~II's central concern.

Part~II asks what a set-membership predicate built from atoms looks like. A set
$S$ over a universe induces a Boolean query per element; an approximate
set\index{approximate set} makes each query a $\Bernoulli{\BoolType}$
instance, mutually independent by Axiom~1 given the realized rates. The canonical construction is the Bloom filter\index{Bloom filter}, a mechanism for
manufacturing one-sided instances ($\fnrate = 0$) with a tunable false-positive
rate, developed in \Cref{ch:approximate-sets}; chapter~6 then propagates error
through set operations, chapter~7 develops the classification measures, and
chapter~8 derives the lower bounds that make space-accuracy duality precise. The
Boolean thread becomes the foundation, one $\Bernoulli{\BoolType}$ instance at a
time.
